\chapter{Drei Fragen über die christliche Kirche}
\quellseite{1}

\section[3. Mit welchen Kirchen sollen wir Gemeinschaft halten?]{3. Mit welchen Kirchen sollen wir Gemeinschaft halten und von welchen sollen wir uns trennen?}

\subsection[a. Gemeinschaft rechtgläubiger Kirchen]{a. Eine rechtgläubige Kirche soll mit andern rechtgläubigen Kirchen Gemeinschaft halten}

Die drei Fragen über die christliche Kirche, die wir hier in unserm „Lutheraner“ beantworten wollten, waren die folgenden: erstens, was eigentlich und wahrhaftig die Kirche sei; zweitens, wie die Kirche hier auf Erden erscheine; drittens, mit welchen Kirchen wir Gemeinschaft halten und von welchen wir uns trennen sollen. Die ersten beiden Fragen haben wir schon im vorigen Jahrgang beantwortet, und wir stehen nun bei der dritten Frage. Und da haben wir auch schon aus Gottes Wort gehört, daß eine \gesperrt{rechtgläubige Kirche mit andern rechtgläubigen Kirchen Gemeinschaft halten soll}. Wir wollen aber diesen wichtigen Punkt noch etwas weiter besehen und deshalb noch ein anderes bekanntes Gotteswort nehmen, welches von anderer Seite her unsern Satz beleuchten wird — unsern Satz, daß eine \gesperrt{rechtgläubige Kirche mit andern rechtgläubigen Kirchen Gemeinschaft halten soll}.

Eph. 4, 1–6. steht also geschrieben: \gesperrt{„So ermahne nun euch ich Gefangener in dem HErrn, daß ihr wandelt, wie sich’s gebührt eurem Berufe, darinnen ihr berufen seid, mit aller Demuth und Sanftmuth, mit Geduld, und vertraget einer den andern in der Liebe, und seid fleißig zu halten die Einigkeit im Geist durch das Band des Friedens. Ein Leib und Ein Geist, wie ihr auch berufen seid auf einerlei Hoffnung eures Berufs. Ein HErr, Ein Glaube, Eine Taufe, Ein Gott und Vater (unser) aller, der da ist über euch alle, und durch euch alle, und in euch allen.“}

Wieder eine apostolische Ermahnung. Wieder eine Ermahnung um Christi willen. Wieder eine überaus ernste und eindringliche Ermahnung.

An wen?

Zunächst an die aus Heiden und Juden gesammelte Kirche in Ephesus, der Hauptstadt Kleinasiens.

Zwischen Heiden und Juden war naturgemäß immer ein großer Abstand. Dieser Abstand hatte sich in Kleinasien und sonderlich in Ephesus zu bitterer Feindschaft verschärft, nachdem der Kaiser Augustus den in Kleinasien wohnenden Juden besondere Rechte und Freiheiten gewährt hatte. Und es war Gefahr, daß etwas von dieser Feindschaft, von diesem Abstand und Gegensatz, sich in die, wie gesagt, aus Juden und Heiden gesammelte Kirche zu Ephesus verschleppen und da Zwietracht und gar offene Spaltung hervorrufen möchte; das um so mehr, als sich bei den Judenchristen das alte Jüdische und bei den Heidenchristen das alte Heidnische doch immer noch — wenn auch bei den aufrichtigen Christen nur als Schwachheit — mehr oder weniger geltend machte und merken ließ. Altüberliefertes und Altgewohntes ist auch nach der Bekehrung nicht so auf einmal verschwunden; der Regel nach gehört dazu die lange und geduldige Arbeit des Heiligen Geistes durchs Wort. Daher ermahnte der Heilige Geist die Christenheit zu Ephesus, daß sie wandeln sollten, wie sich’s gebührt ihrem Berufe, als rechte Jünger Jesu; mit aller Demuth und Sanftmuth sollten sie wandeln, und mit Geduld; und es sollte einer den andern vertragen in der Liebe; und sie sollten fleißig sein zu halten die Einigkeit im Geist durch das Band des Friedens; sie sollten bedenken, daß sie, einerlei ob sie aus Juden oder aus Heiden kommen, nun Ein geistlicher Leib seien in Christo und von Einem Heiligen Geiste regiert werden und Eine Hoffnung der Seligkeit haben, und daß sie Einen HErrn haben und Einen Glauben und Eine Taufe und Einen Gott und Vater, der über sie alle herrsche und durch sie alle wirke und in ihnen allen wohne.

Dieselbe Ermahnung, mit all ihrem Ernste und mit all ihrer Eindringlichkeit, gilt auch den rechtgläubigen Kirchen unserer Tage.

\quellseite{2}
Die einzelnen jetzt rechtgläubigen Kirchen und kirchlichen Verbindungen haben eine verschiedene Vergangenheit und Geschichte. Die einen stehen schon seit langer Zeit unter dem Panier des lauteren Bekenntnisses und unter dem heilsamen Einfluß der unverfälschten Lehre und im Zusammenhang mit andern rechtgläubigen Kirchen. Was sie früher und ursprünglich etwa waren, das ist vergessen, das zeigt sich nicht mehr, das übt keinen Einfluß mehr auf sie aus. — Die andern stehen erst seit kürzerer Zeit unter dem Panier des lauteren Bekenntnisses und unter dem heilsamen Einfluß der unverfälschten Lehre und im Zusammenhang mit andern rechtgläubigen Kirchen. Sie waren früher und ursprünglich etwa unirt, oder doch nicht in Wahrheit und in allen Stücken lutherisch, rechtgläubig. Jetzt sind sie rechtgläubig und wollen es in Wahrheit sein. Aber es geht ihnen wie einem Menschen, der eben von einer langen Krankheit genesen ist: die Nachwehen des alten krankhaften Zustandes sind noch da, einzelne Glieder sind noch recht schwach, hie und da droht ein Rückfall. Es fehlt ihnen wohl auch an der rechten Erfahrung, an der rechten Selbsterkenntniß, an der rechten Vorsicht. Sie halten sich für stärker, als sie in der That sind.

Mit diesem beispielsweise angeführten zweifachen Unterschied unter den rechtgläubigen Kirchen wollen wir uns hier begnügen. Wir wollen der Kürze und Bequemlichkeit wegen die ersteren „die Alten“ und die letzteren „die Neuen“ nennen.

Da ist nun Gefahr, daß zwischen diesen Uneinigkeit entsteht. Wie?

Die Alten sehen etwa mehr auf die alte Krankheit und deren Nachwehen als auf die neue Genesung und deren feine Erscheinungen bei den Neuen. Und sie verlangen jetzt Vollkraft. Und sie wollen bessern. Nun, das sollen sie auch. Aber sie fassen die Sache in unrichtiger, in fleischlicher Weise an: sie gebärden sich als Vormünder, sie machen Anspruch auf unbedingten und sofortigen Gehorsam; mit trauriger Stirne sehen sie die Neuen an, sie arbeiten sich in einen richterlichen und abgünstigen Sinn hinein; sie kritisiren die Neuen hinter deren Rücken, sie machen ihnen auch feierliche Vorstellungen über ihre Fehler — item, es fehlt ihnen in ihrem Verhalten den Neuen gegenüber an der rechten christlichen Demuth, Sanftmuth, Geduld, Verträglichkeit, in Summa, an der Liebe, ohne welche kein Segen auf ihrem Thun ruhen kann.

Und die Neuen? Die blicken etwa mehr auf ihre neue Genesung und deren erfreuliche Symptome als auf ihre alte lange Krankheit und deren noch vorhandene Nachwehen. Sie meinen, sie haben jetzt Vollkraft. Sie wollen sich von den Alten nicht meistern lassen. Das nöthige Bessern wollen sie selber besorgen. Anstatt das, was ihnen von den Alten gesagt wird, ganz sachlich anzusehen, machen sie ihre Augen weit auf gegen die etwa verkehrte Weise, in der die Alten das Ding angreifen, kritisiren nun auch ihrerseits an den Alten herum und — weisen Gutes und Böses, Göttliches und Menschliches gleicher Weise zurück. Auch ihnen, in ihrem Verhalten den Alten gegenüber, fehlt es an der rechten christlichen Demuth, Sanftmuth, Geduld, Verträglichkeit, in Summa, an der Liebe, ohne welche auch sie keinen Segen haben werden.

So ist Gefahr, daß zwischen den Alten und den Neuen Uneinigkeit kommt. Ja, so kommt Uneinigkeit. So wird die Einigkeit im Geist, im Glauben, im Bekenntniß, in der Lehre, gefährdet und endlich zerstört. — Und das ist so nicht nur bei den „Alten“ und den „Neuen“, die ich beispielsweise angeführt habe, sondern das ist überhaupt ähnlich so unter den rechtgläubigen Kirchen. Durch Mangel an christlicher Demuth, Sanftmuth, Geduld, Verträglichkeit, item, an Liebe, wird die Einigkeit im Geist gefährdet und zerstört.

Da mahnt nun der heilige Geist in unserm Gottesworte die rechtgläubigen Kirchen, daß sie ihres Christenberufes eingedenk seien und demgemäß unter einander handeln und wandeln. Hilf Gott! wahrhaftig! wie einzelne Christen unter einander handeln und wandeln sollen, so sollen auch ganze Kirchen unter einander handeln und wandeln. Und nun gar solche Kirchen, denen Gott die Gnade der Rechtgläubigkeit aus Gnaden geschenkt hat und die er ganz sonderlich zum Licht der Welt gemacht hat!

Und wie sollen diese demgemäß unter einander handeln und wandeln?

In Liebe.

Ja, in Liebe. Rechtgläubige Kirchen sollen, um ihre Glaubens- und Bekenntnißeinigkeit zu erhalten, in Liebe unter einander handeln und wandeln. Vermuthe mir hier niemand Unionisterei! Citire mir hier niemand Worte Luthers, der da sagt, daß die Liebe in Glaubenssachen nicht mitzureden habe, der aber nur meint, daß die Liebe nicht wider den Glauben gelten dürfe. Lies 1 Cor. 13 und citire dir selbst da V. 2.: „Wenn ich weissagen könnte, und wüßte alle Geheimnisse, und alle Erkenntniß, und hätte allen Glauben, also, daß ich Berge versetzte, und hätte der Liebe nicht; so wäre ich nichts.“

Die Liebe soll wogen und wallen von den Alten auf die Neuen und von den Neuen auf die Alten, von denen, die bessern wollen, auf die, die gebessert werden sollen, und von diesen auf jene.

Und bei beiden soll die Liebe erscheinen in Demuth, in aller Demuth, in gründlicher, wahrhaftiger, sich durch alles durchziehender Demuth, und in folgender Sanftmuth und in wiederum folgender Geduld und Langmuth und darin, daß die einen die andern vertragen, eben in der Liebe. — Ja, solange der alte Adam noch nicht endgültig todt ist, hat auch in den rechtgläubigen Kirchen die Liebe viel Arbeit, sich zu bewähren in aller Demuth und Sanftmuth und Geduld und Verträglichkeit.

Aber die Liebe soll sich so bewähren — und zwar zuerst und vor allen Dingen bei den Alten, bei denen, die andere bessern wollen. Sonst werden diese einen großen Fall thun und narren und schaden und kein nütze sein.

Und so, in dieser fein mit aller Demuth und Sanftmuth und Geduld und Verträglichkeit einhergehenden Liebe sollen rechtgläubige Kirchen \gesperrt{fleißig sein zu halten die Einigkeit im Geist durch das Band des Friedens}.

Dazu hat sie der Herr berufen. Das liegt in ihrem Christenberuf. So will Gott, daß seine rechtgläubigen Kirchen unter einander sein und thun sollen. Sie sollen fleißig sein, sie sollen einen Eifer, einen göttlichen Eifer haben, sie sollen alle Mühe, alle selbstverleugnende, alle alles tragende Mühe anwenden — wozu? Dazu, daß sie unter einander halten, erhalten, festhalten die Einigkeit im Geist, die Glaubenseinigkeit, die Bekenntnißeinigkeit, die Einigkeit in einerlei Rede, die Einigkeit in einem Sinn und in einerlei Meinung; und daß alsdann der Friede, als ein sie fest und innig umschlingendes Band, sie unter einander und denen gegenüber, die draußen sind, verbinde.

\quellseite{3}
Wie ist es nun? Wie soll sich hiernach die Sache in der Praxis stellen?

Es kränkelt etwa eine rechtgläubige Kirche. Man nimmt krankhaftes an ihr wahr, in ihrem Glauben, in ihrer Liebe, in ihrer Lehrweise, in ihren Aeußerungen, in ihrem Thun und Treiben. Man ermahnt sie deshalb. Sie stellt sich ungebärdig. — Das alles ist nicht recht. — \gesperrt{Aber solange eine solche Kirche nicht offenbar und klarlich falsche Lehre führt oder in ihrer Mitte duldet; solange eine Kirche falsche Lehre nicht hartnäckig vertheidigt oder sich entschieden weigert, sie aus ihrer Mitte zu verbannen: so lange soll man sie nicht als eine falschgläubige halten und behandeln und verlassen, sondern so lange soll man fleißig sein zu halten die Einigkeit im Geist durch das Band des Friedens}, in der Weise, die unser Gotteswort anzeigt.

Verstößt auch eine Mutter ihr Kind, wenn es kränkelt und sich dabei ungebärdig stellt? Sie wäre keine Mutter mehr, sie wäre des Namens unwerth, sie hätte all ihren Beruf verleugnet, wenn sie das thäte. — Welche Kirche nun eine andere Kirche zu ermahnen hat und also Mutterstelle vertreten will, die merke sich dies. Stößt auch Christus, unser Herr, so schnell den Leuchter einer Kirche von seiner Stätte?

Und wenn nun eine so kränkelnde und sich bislang ungebärdig stellende Kirche sich ändert und sich willig erzeigt, die Ermahnung anzunehmen und ihr Folge zu geben, was dann? \gesperrt{Dann soll man erst recht und desto mehr Liebe an ihr beweisen, sie aufs freundlichste und innigste in die Gemeinschaft ziehen, aufs allerliebreichste und mit dem größten Vertrauen sie aufnehmen} — „auf daß wir nicht übervortheilt werden vom Satan. Denn uns ist nicht unbewußt, was er im Sinn hat“, 2 Cor. 2, 11.

Sieht nicht jeder Christ ein, daß Gott haben will, daß Christen so handeln sollen? Oder meint jemand, daß Gott eine sonderliche Kirchenpolitik will? daß er für die Kirchen ein anderes Verhalten vorschreibt als für die einzelnen Christen? Wo das?

Oder — wir wollen jetzt die andere Seite betrachten — es hat eine rechtgläubige Kirche in einem gewissen Eifer um die reine Lehre eine andere, die Anlaß zu Verdacht und Tadel gegeben hat, verkehrt, unbrüderlich, lieblos behandelt und sie dadurch schwer gekränkt und ihr großen Anstoß gegeben; was soll diese andere Kirche dann thun? Soll sie sich dann auf den Richterstuhl setzen, aufprotzen, verdammen und zum wenigsten Buße in Sack und Asche verlangen?

Nein, sie soll die Liebe walten lassen. Sie soll wandeln und handeln, wie sich’s gebührt ihrem Berufe, darinnen sie berufen ist; mit aller Demuth und Sanftmuth, mit Geduld; sie soll vertragen in der Liebe; \gesperrt{sie soll fleißig sein zu halten die Einigkeit im Geist durch das Band des Friedens}. Sie soll in aller Demuth anerkennen, daß sie doch Anlaß gegeben hat zu Verdacht und Tadel. Sie soll mit Sanftmuth und Geduld in ungeschwächter Liebe das angethane Unrecht vertragen. \gesperrt{Sie soll fleißig sein zu halten die Einigkeit im Geist durch das unmißverständlichste Bekenntniß zur Wahrheit und durch den ausdrücklichen Widerruf aller mißverständlichen oder verkehrten Reden}. Sie soll alles, alles thun, um das gelockerte oder gar gelöste Band des Friedens wieder zu festigen, zu knüpfen. Und all das soll sie mit desto größerer Macht und Herzlichkeit thun, wenn sie sieht, daß ihre Schwesterkirche auch nur beginnt einzusehen, daß sie zu weit gegangen ist.

Und dann soll wiederum auch diese sich nicht schämen, ihr Unrecht voll und ganz mit aller Demuth einzugestehen und, soweit es noch irgend möglich ist, wieder gut zu machen.

\gesperrt{„Ein Leib und Ein Geist, wie ihr auch berufen seid auf einerlei Hoffnung eures Berufs. Ein HErr, Ein Glaube, Eine Taufe, Ein Gott und Vater unser aller, der da ist über euch alle, und durch euch alle, und in euch allen.“} Mit diesen Worten hält der Heilige Geist der rechtgläubigen Christenheit ein so köstliches und Gott so wohlgefälliges Ding vor, daß alle rechtgläubigen Kirchen mit dem größten Eifer und Fleiß und mit sich selbst und ihre Privatinteressen und -Stimmungen völlig verleugnender Mühe diese Einigkeit im Geist durch das Band des Friedens zu bewahren suchen sollen, wo sie da ist, und zu halten suchen sollen, wo sie weichen will.

Welche rechtgläubige Kirche das nicht thut, die macht sich selbst zum Aas, um das die Adler sich sammeln werden.

C. M. Z.\par
